\section{Natural mass and the critical relative-cancellation threshold}
\label{sec:natural-mass-threshold}

\begin{definition}[Natural post-bin mass regime]
\label{def:natural-mass-regime}
The literal coefficient is in the natural post-bin mass regime if
\begin{equation}
 M_X\ge X^{-o(1)}N_{0,X}.
\label{eq:natural-mass-lower}
\end{equation}
Together with \eqref{eq:mass-envelope}, this means
\[
 M_X=X^{o(1)}N_{0,X},
\]
where the notation denotes two-sided comparability up to subpower
factors.
\end{definition}

TPC-83 proves that \eqref{eq:natural-mass-lower} is equivalent to
natural determinant-energy survival:
\begin{equation}
 D_X\ge
 X^{-o(1)}\frac{N_{0,X}^2}{Q_X}.
\label{eq:natural-energy-lower}
\end{equation}
Combined with \eqref{eq:energy-envelope}, this gives two-sided
subpower comparability for \(D_X\).

\begin{theorem}[Natural mass forces natural effective support]
\label{thm:natural-effective-support}
In the natural post-bin mass regime,
\begin{equation}
 \boxed{
 \SeffX=X^{o(1)}Q_X.}
\label{eq:natural-effective-support}
\end{equation}
Equivalently,
\[
 X^{-o(1)}Q_X\le\SeffX\le X^{o(1)}Q_X.
\]
Consequently,
\begin{equation}
 S_X\ge\SeffX\ge X^{-o(1)}Q_X,
\label{eq:broad-actual-support}
\end{equation}
recovering the broad-actual-support conclusion of TPC-83.
\end{theorem}

\begin{proof}
Use the mass lower bound and the energy upper bound:
\[
 \SeffX=\frac{M_X^2}{D_X}
 \ge
 \frac{
 X^{-o(1)}N_{0,X}^2
 }{
 X^{o(1)}N_{0,X}^2/Q_X
 }
 =
 X^{-o(1)}Q_X.
\]
For the reverse inequality, use the mass upper bound and the
TPC-83 energy lower bound:
\[
 \SeffX
 \le
 \frac{
 X^{o(1)}N_{0,X}^2
 }{
 X^{-o(1)}N_{0,X}^2/Q_X
 }
 =
 X^{o(1)}Q_X.
\]
Finally, \(\SeffX\le S_X\) by
\cref{thm:exact-factorization}.
\end{proof}

\begin{remark}[No phase information has been gained]
\Cref{thm:natural-effective-support} fixes the magnitude-spreading
level up to subpower factors.  It gives no upper bound for
\(\canc_X\).  The common-phase and balanced-phase examples in
\cref{ex:same-magnitude-opposite-phase} both have maximal effective
support.
\end{remark}

\subsection{A general threshold conversion}

\begin{theorem}[Effective-support to cancellation conversion]
\label{thm:threshold-conversion}
Suppose, for fixed exponents \(\sigma,\chi\ge0\), that
\[
 \SeffX=X^{\sigma+o(1)}.
\]
Then
\begin{equation}
 \boxed{
 \Flat_X\le X^{\chi+o(1)}
 \quad\Longleftrightarrow\quad
 \canc_X\le X^{(\chi-\sigma)/2+o(1)}.}
\label{eq:general-threshold-conversion}
\end{equation}
When \(\sigma>\chi\), the right side is a genuine relative
cancellation requirement.
\end{theorem}

\begin{proof}
The exact factorization gives
\[
 \canc_X^2=\frac{\Flat_X}{\SeffX}.
\]
Insert the two-sided asymptotic for \(\SeffX\) in either direction
and take square roots.
\end{proof}

\begin{theorem}[Critical \(1/J\) equivalence]
\label{thm:critical-one-over-J}
Assume the natural post-bin mass regime at
\[
 \beta=\frac{267}{400}.
\]
Then the TPC-32 flatness threshold and the relative cancellation
requirement are equivalent:
\begin{equation}
 \boxed{
 \begin{aligned}
 \Flat_X\le X^{1/400+o(1)}
 &\quad\Longleftrightarrow\quad
 \canc_X=\frac{|Z_X|}{M_X}\\
 &\le
 X^{-133/400+o(1)}
 =
 \frac{X^{o(1)}}{J_X}.
 \end{aligned}}
\label{eq:critical-one-over-J}
\end{equation}
\end{theorem}

\begin{proof}
By \cref{thm:natural-effective-support},
\[
 \SeffX=X^{267/400+o(1)}.
\]
Apply \cref{thm:threshold-conversion} with
\(\sigma=267/400\) and \(\chi=1/400\):
\[
 \frac{\sigma-\chi}{2}
 =
 \frac12\left(\frac{267}{400}-\frac1{400}\right)
 =
 \frac{133}{400}.
\]
The critical scale relation gives
\[
 J_X=X^{133/400+o(1)}.
\]
\end{proof}

\begin{corollary}[Relative and absolute zero-mode targets]
\label{cor:relative-absolute}
In the same natural mass regime, the following three statements are
equivalent up to subpower factors:
\begin{align}
 \Flat_X&\le X^{1/400+o(1)},
 \label{eq:equiv-flatness}\\
 \frac{|Z_X|}{M_X}&\le\frac{X^{o(1)}}{J_X},
 \label{eq:equiv-relative}\\
 |Z_X|&\le X^{o(1)}Q_X^2.
 \label{eq:equiv-absolute}
\end{align}
\end{corollary}

\begin{proof}
The first two are equivalent by
\cref{thm:critical-one-over-J}.  Natural mass gives
\[
 X^{-o(1)}J_XQ_X^2
 \le M_X\le
 X^{o(1)}J_XQ_X^2.
\]
Multiplying \eqref{eq:equiv-relative} by the upper bound for \(M_X\)
gives \eqref{eq:equiv-absolute}.  Conversely, divide
\eqref{eq:equiv-absolute} by the lower bound for \(M_X\).
\end{proof}

\begin{remark}[Why the factor is \(1/J\)]
The \(1/J_X\) is not inserted heuristically.  It is the square-root
gap between the natural effective-support exponent \(267/400\) and
the allowed flatness exponent \(1/400\):
\[
 \frac12\left(\frac{267}{400}-\frac1{400}\right)
 =
 \frac{133}{400}.
\]
\end{remark}
